\chapter{REVISÃO DA LITERATURA}

O uso de \textit{smartphones} tornou-se parte constante da rotina contemporânea. O dispositivo concentra atividades de comunicação, acesso à informação, redes sociais, entretenimento e trabalho. Entretanto, a literatura recente indica que o uso excessivo pode estar associado a impactos relevantes na saúde física, mental, cognitiva e social.

Este capítulo apresenta o cenário atual de uso, detalha os principais efeitos identificados em estudos científicos, analisa o papel do \textit{design persuasivo} no aumento do engajamento e discute as limitações das ferramentas existentes para controle do tempo de tela.

\section{Cenário atual e o uso de \textit{smartphones}}

O Brasil está entre os países com elevado tempo médio de exposição a telas digitais. Segundo dados divulgados pelo Jornal da USP (2023), brasileiros passam em média cerca de nove horas por dia diante de telas, considerando \textit{smartphones} e computadores. Já a Folha de S.Paulo (2024) indica que, ao considerar apenas \textit{smartphones}, o tempo médio diário é de aproximadamente três horas e meia. Esses dados evidenciam um padrão de uso intenso e reforçam a necessidade de analisar os possíveis impactos associados ao tempo prolongado de tela.

\section{Impactos do uso excessivo de \textit{smartphones}}

Nesta seção são apresentados os principais impactos associados ao uso excessivo de \textit{smartphones}, organizados em quatro áreas: física, cognitiva, emocional (incluindo efeitos sobre o sono) e social/comportamental. O objetivo é sintetizar evidências da literatura que indiquem de que forma o tempo prolongado de tela pode afetar o corpo, os processos mentais, o bem-estar psicológico e as relações interpessoais dos usuários.

\subsection{Impactos na saúde física}

Entre os efeitos físicos mais relatados está o fenômeno conhecido como "\textit{text neck}". O estudo de Neupane et al. (2017) descreve que a inclinação prolongada da cabeça durante o uso do \textit{smartphone} está associada a dores cervicais, dores de cabeça e desconforto nos ombros. Esse padrão postural repetitivo pode comprometer o bem-estar e a qualidade de vida quando mantido por longos períodos. Além disso, o tempo excessivo de tela está associado à redução da atividade física, contribuindo para comportamentos sedentários.

\subsection{Impactos cognitivos}

O uso intenso do \textit{smartphone} também pode interferir na capacidade de concentração. Ward et al. (2017) demonstraram que a simples presença do aparelho, mesmo quando não está em uso, pode reduzir o desempenho em tarefas que exigem atenção sustentada. A exposição constante a notificações e a alternância frequente entre aplicativos podem dificultar o foco contínuo e aumentar a dispersão. Wilmer et al. (2017) associam o uso excessivo à redução da memória de trabalho e ao prejuízo do controle inibitório, o que pode impactar o desempenho acadêmico e profissional.

\subsection{Impactos emocionais e no sono}

Pesquisas indicam associação entre uso excessivo de \textit{smartphones} e aumento de sintomas de ansiedade, depressão e distúrbios do sono \cite{twenge2018screentime}. O uso prolongado, especialmente em redes sociais, pode estimular comparações sociais frequentes, sensação de inadequação e estresse.

Twenge e Campbell (2018) observaram, nos Estados Unidos, crescimento de sintomas depressivos entre jovens com maior tempo de exposição a telas. Esses achados reforçam a necessidade de compreender a relação entre padrões de uso e saúde mental.

Outro ponto relevante é o impacto no sono. Chang et al. (2015) demonstraram que a exposição à luz azul das telas antes de dormir interfere na produção de melatonina, hormônio responsável pela regulação do sono. Esse efeito pode atrasar o início do sono e reduzir sua qualidade, afetando memória, humor e concentração ao longo do tempo.

\subsection{Impactos sociais e comportamentais}

O uso excessivo também pode afetar relações interpessoais. Roberts e David (2016) introduzem o conceito de "\textit{phubbing}", caracterizado por priorizar o \textit{smartphone} em detrimento de interações presenciais. Esse comportamento está associado à redução da qualidade das relações e aumento da sensação de isolamento. Outro conceito apresentado é o de "\textit{continuous partial attention}", no qual o indivíduo mantém a atenção dividida entre várias atividades ao mesmo tempo, um dos fatores principais que colaboram para isso é a necessidade constante de responder notificações. Esse padrão pode gerar sensação de urgência contínua e níveis mais elevados de estresse.

\section{\textit{Design persuasivo} e engajamento do usuário}

O aumento do tempo de uso não depende apenas de hábitos individuais. Muitos aplicativos são projetados com estratégias conhecidas como \textit{design persuasivo}, cujo objetivo é maximizar retenção e engajamento do usuário. Essas estratégias incluem notificações frequentes, recompensas variáveis, personalização de conteúdo e mecanismos de comparação social \cite{chen2023}. Tais elementos estimulam verificações recorrentes e prolongam o tempo de permanência na plataforma. Ao compreender que esses mecanismos são intencionalmente utilizados para aumentar o engajamento, torna-se relevante discutir alternativas que promovam o uso mais consciente da tecnologia.

\section{Ferramentas atuais de controle de tempo de tela e suas limitações}

Como resposta ao aumento do tempo de uso, fabricantes passaram a oferecer ferramentas nativas de monitoramento e controle. No \textit{Android}, o \textit{Digital Wellbeing} permite visualizar relatórios de uso, definir limites diários para aplicativos e ativar modos de foco. No \textit{iOS}, o \textit{Screen Time} oferece funcionalidades semelhantes, incluindo bloqueio por senha e gerenciamento familiar.

Além das opções nativas, existem aplicativos de terceiros como \textit{RescueTime}\footnote{RescueTime. Disponível em: \url{https://www.rescuetime.com/}. Acesso em: 18 fev. 2026.}, \textit{Freedom}\footnote{Freedom. Disponível em: \url{https://freedom.to/}. Acesso em: 18 fev. 2026.} e \textit{Forest}\footnote{Forest. Disponível em: \url{https://www.forestapp.cc/}. Acesso em: 18 fev. 2026.}. O \textit{RescueTime} monitora e gera relatórios detalhados de uso. O \textit{Freedom} permite bloqueio programado de aplicativos. O \textit{Forest} utiliza abordagem gamificada, incentivando períodos de foco por meio de recompensas simbólicas.

Apesar dessas alternativas, estudos indicam que tais recursos podem ser facilmente contornados pelos próprios usuários \cite{chen2023}. Limites podem ser estendidos, configurações alteradas e aplicativos desinstalados. Dessa forma, a efetividade dessas soluções depende fortemente do autocontrole individual, o que representa uma limitação importante.

\section{Discussão}

A literatura demonstra que o uso excessivo de \textit{smartphones} está associado a impactos físicos, cognitivos, emocionais e sociais. Além disso, muitos aplicativos são desenvolvidos com estratégias que incentivam acessos frequentes e permanência prolongada, o que contribui para padrões de uso mais intensos ao longo do tempo.

No campo cognitivo, as evidências indicam que o problema não se limita apenas ao tempo de uso ativo. Ward et al. (2017) demonstram que a simples presença do \textit{smartphone} já pode prejudicar tarefas que exigem concentração. Isso sugere que a relação com o dispositivo vai além do uso em si e envolve também a constante expectativa de interação. Assim, reduzir apenas o tempo de tela pode não ser suficiente, é necessário oferecer estratégias que ajudem o usuário a reorganizar sua rotina e sua relação com o aparelho.

Pode parecer contraditório propor um aplicativo para reduzir o tempo de uso do próprio \textit{smartphone}, já que isso envolve a utilização do mesmo dispositivo. No entanto, essa escolha é intencional. O público-alvo deste trabalho é composto principalmente por jovens, que já utilizam o celular como principal meio de comunicação, estudo e entretenimento. Nesse contexto, o próprio aparelho pode ser utilizado como ferramenta de apoio para promover mudanças de hábito.

A proposta não é criar mais um aplicativo que dispute a atenção do usuário, mas sim uma ferramenta objetiva e de uso pontual, voltada ao acompanhamento de metas e ao monitoramento do tempo de tela. O foco está em fornecer informações claras, estimular a definição de limites e acompanhar o progresso ao longo do tempo. Dessa forma, o aplicativo atua como um meio para reduzir o uso excessivo, e não para ampliá-lo.

Outro ponto central da proposta é o aspecto coletivo. Mudanças de hábito tendem a ser mais sustentáveis quando ocorrem com apoio social. Ao permitir que usuários participem de desafios em grupo, comparem resultados e incentivem uns aos outros, o aplicativo utiliza a dinâmica coletiva como mecanismo de motivação. Estratégias semelhantes já são utilizadas com sucesso em aplicativos de treino, como o \textit{Gym Rats}\footnote{Gym Rats. Disponível em: \url{https://www.gymrats.com.br/}. Acesso em: 18 fev. 2026.}, que organiza desafios de treinos focados em metas de desempenho e consistência ao longo de períodos determinados, combinando rotina estruturada, acompanhamento em grupo e reforços motivacionais. De forma análoga, aplicativos de aprendizado de idiomas como o \textit{Duolingo}\footnote{Duolingo. Disponível em: \url{https://www.duolingo.com/}. Acesso em: 18 fev. 2026.} utilizam lições curtas, sistema de pontos, metas diárias e comparação entre usuários para estimular a prática contínua. Nesses contextos, a gamificação e a dinâmica social têm se mostrado eficientes para aumentar a adesão e a permanência dos participantes. Aplicar essa lógica à redução do tempo de tela pode tornar o processo mais motivador e contínuo.

Dessa forma, a literatura aponta para a necessidade de soluções que vão além do simples monitoramento ou bloqueio técnico. É nesse contexto que se justifica o desenvolvimento de um aplicativo gamificado em grupo, que utilize o próprio \textit{smartphone} como ferramenta de apoio para promover um uso mais equilibrado e consciente da tecnologia.